\chapter{Mô hình và kỹ thuật nền tảng}

\section{Kiến trúc web API với FastAPI}

Backend JobConnect dùng FastAPI làm lớp API runtime. Đây không phải lựa chọn được ghi chung chung là ``web framework'', mà gắn trực tiếp với dependency \texttt{fastapi==0.115.6} trong \texttt{backend/requirements.txt}. FastAPI phù hợp với dự án vì cung cấp request validation, response model, dependency injection, OpenAPI generation và khả năng tổ chức router theo module. Trong repository, app entrypoint nằm tại \texttt{backend/src/jobconnect/main.py}, metadata và router registration nằm tại \texttt{backend/src/jobconnect/app.py}, còn production API router aggregator nằm tại \texttt{backend/src/jobconnect/modules/api/router.py}.

Theo tài liệu kiến trúc \cite{hld-architecture}, runtime được chia thành các lớp: API layer, auth/authorization layer, domain services, data access layer, worker layer và external services. Cách chia này tránh để endpoint chứa toàn bộ logic nghiệp vụ. API layer chịu trách nhiệm validate request và trả schema rõ ràng; service layer xử lý trạng thái nghiệp vụ như activate resume, publish job, accept invite; data layer làm việc với PostgreSQL; worker layer xử lý parse/embedding; external services được bọc qua integration boundary.

\begin{figure}[H]
\centering
\fbox{\begin{minipage}{0.88\textwidth}
\centering
\textbf{FastAPI app}\\[0.3em]
auth -- users/profiles -- organizations -- documents -- resumes -- jobs\\
matching/search -- applications -- invites -- notifications -- admin\\[0.8em]
\textbf{Workers}\\
object storage read -> extraction -> preprocessing -> parser -> embeddings\\[0.8em]
\textbf{PostgreSQL + pgvector}\\
business entities + resume/job embedding tables
\end{minipage}}
\caption{Kiến trúc runtime tổng quát của JobConnect}
\docsrc{\cite{hld-architecture}.}
\end{figure}

FastAPI còn quan trọng vì OpenAPI là contract surface cho frontend. Tài liệu API LLD yêu cầu mọi route có request/response schema explicit, mọi enum phải thể hiện trong OpenAPI và mọi lỗi business phải đi qua error envelope thống nhất \cite{lld-api}. Vì vậy hệ thống không nên trả response tùy tiện theo từng endpoint; nó cần schema rõ cho \texttt{AuthResponse}, \texttt{CandidateResumeSummary}, \texttt{JobPostSummary}, \texttt{ApplicationSummary}, \texttt{InviteSummary}, semantic search responses và matching responses.

\section{Pydantic, schema validation và error envelope}

Pydantic \texttt{2.11.7} được dùng cho request/response model. Vai trò của Pydantic trong hệ thống không chỉ là kiểu dữ liệu Python, mà là lớp bảo vệ API contract. Các request như create resume, update job, semantic search, application status update hoặc invite request phải reject field lạ nếu endpoint không cho phép metadata JSON. Theo implementation matrix, Pydantic models dùng \texttt{extra="forbid"} ở nhiều nơi, và PATCH endpoints đã được điều chỉnh theo hướng partial update bằng \texttt{exclude\_unset} \cite{implementation-matrix}.

Error envelope target có dạng:

\begin{lstlisting}[caption={Error envelope target trong API contract}]
{
  "error": {
    "code": "validation_error",
    "message": "Request validation failed.",
    "fields": {
      "email": "Invalid email address."
    },
    "request_id": "req_01hxyz"
  }
}
\end{lstlisting}

Cách chuẩn hóa lỗi giúp frontend xử lý thống nhất các trường hợp \texttt{400}, \texttt{401}, \texttt{403}, \texttt{404}, \texttt{409}, \texttt{413}, \texttt{415}, \texttt{422} và \texttt{500}. Với workflow upload, \texttt{413} dùng cho file vượt giới hạn, \texttt{415} cho MIME type không hỗ trợ. Với lifecycle, \texttt{409} phù hợp cho duplicate invite, duplicate application hoặc invalid transition.

\section{PostgreSQL và pgvector}

PostgreSQL là source of truth cho business entities: \texttt{users}, \texttt{organizations}, \texttt{candidate\_profiles}, \texttt{recruiter\_profiles}, \texttt{candidate\_resumes}, \texttt{job\_posts}, \texttt{uploaded\_documents}, \texttt{parse\_jobs}, \texttt{applications}, \texttt{application\_events}, \texttt{recruiter\_invites}, \texttt{notifications} và \texttt{audit\_logs}. pgvector được dùng trong cùng PostgreSQL để lưu semantic vectors cho resume/job fields. Điều này giúp hệ thống giữ transactional business data và semantic retrieval data trong một data store quản lý được ở MVP.

Schema LLD quy định primary key dùng \texttt{BIGINT GENERATED BY DEFAULT AS IDENTITY}; timestamps dùng \texttt{TIMESTAMPTZ}; enum-like values dùng \texttt{TEXT} hoặc \texttt{VARCHAR} với \texttt{CHECK} constraints thay vì PostgreSQL native enum; normalized skills/certifications dùng \texttt{TEXT[]} với default \texttt{\{\}}; vector columns dùng \texttt{VECTOR(384)} và nullable \cite{lld-schema}. Việc để vector nullable là thiết kế có chủ đích: thiếu embedding không làm fail toàn bộ matching run, mà component tương ứng được score 0.

\begin{table}[H]
\caption{Nhóm bảng chính trong schema JobConnect}
\centering
\begin{tabularx}{\textwidth}{p{0.30\textwidth}X}
\toprule
\textbf{Nhóm} & \textbf{Bảng cụ thể}\\
\midrule
Identity và profile & \texttt{users}, \texttt{candidate\_profiles}, \texttt{recruiter\_profiles}, \texttt{organizations}\\
Marketplace entities & \texttt{candidate\_resumes}, \texttt{job\_posts}\\
Document pipeline & \texttt{uploaded\_documents}, \texttt{parse\_jobs}\\
Semantic vectors & \texttt{candidate\_resume\_embeddings}, \texttt{job\_post\_embeddings}\\
ATS-lite & \texttt{applications}, \texttt{application\_events}, \texttt{recruiter\_invites}\\
Operations & \texttt{notifications}, \texttt{audit\_logs}\\
\bottomrule
\end{tabularx}
\end{table}

\section{Pipeline xử lý tài liệu CV/JD}

Pipeline ingestion là một trong những phần quan trọng nhất của JobConnect. Theo HLD ingestion, luồng core là: upload file, store original in object storage, create \texttt{uploaded\_documents}, create \texttt{parse\_jobs}, worker marks processing, extract raw text, preprocess text, normalize skills, LLM structured parse, update resume/job fields, expose review payload, save user corrections, generate embeddings, mark succeeded or failed \cite{hld-ingestion}.

\begin{figure}[H]
\centering
\fbox{\begin{minipage}{0.92\textwidth}
\centering
upload file -> object storage -> uploaded\_documents -> parse\_jobs\\
-> extract text -> preprocess -> skill normalization -> LLM/local parser\\
-> draft resume/job fields -> user review -> embeddings -> active/published pool
\end{minipage}}
\caption{Pipeline xử lý CV/JD từ file đến public pool}
\docsrc{\cite{hld-ingestion}.}
\end{figure}

File handling có một invariant quan trọng: parse failure không được xóa original file. Database chỉ lưu metadata và object key/file URL, còn object storage là source of truth cho file gốc. Upload endpoint trả về sau khi file persistence và parse job creation hoàn tất; parse chạy async. Điều này giúp API upload có latency kiểm soát được và tránh giữ kết nối quá lâu trong khi extraction/LLM/embedding chạy.

\subsection{Preprocessing}

Preprocessing phải chạy sau raw text extraction và trước structured parsing. Minimum behavior gồm remove null bytes, invalid control characters, normalize Unicode to NFC cho tiếng Việt, collapse excessive whitespace, preserve technical terms/certification names và không dịch toàn bộ document. Quy tắc ``không dịch toàn bộ document'' rất quan trọng: nếu dịch trước khi parse, thuật ngữ như ``Spring'' có thể bị hiểu sai thành mùa xuân, hoặc certification name bị đổi nghĩa.

\subsection{Skill normalization}

Skill normalization xử lý alias trước hard-filter extraction và embedding. Ví dụ \texttt{reactjs}, \texttt{react.js}, \texttt{React} được normalize thành \texttt{react}; \texttt{nodejs}, \texttt{node.js} thành \texttt{nodejs}; \texttt{pytorch}, \texttt{PyTorch} thành \texttt{pytorch}. Tài liệu yêu cầu gợi ý dictionary MVP khoảng 50-100 high-impact entries \cite{requirements}. Trong code layout hiện tại, phần này tương ứng với \texttt{backend/src/jobconnect/modules/documents/skill\_normalizer.py}.

\subsection{LLM structured parsing}

Parser boundary nằm tại \texttt{backend/src/jobconnect/integrations/llm/}. \texttt{base.py} định nghĩa \texttt{LLMParser} protocol với \texttt{parse\_resume}, \texttt{parse\_job} và \texttt{parser\_version}. \texttt{local.py} là deterministic keyword parser với \texttt{parser\_version = "local-v1"}. \texttt{openai.py} gọi OpenAI-compatible chat completions endpoint với \texttt{response\_format=json\_object}, sau đó validate enum field trước khi trả về. Environment variables gồm \texttt{LLM\_PROVIDER}, \texttt{OPENAI\_API\_KEY}, \texttt{OPENAI\_MODEL}, \texttt{OPENAI\_BASE\_URL}, \texttt{OPENAI\_TIMEOUT\_SECONDS}. Default model trong docs là \texttt{gpt-4o-mini} \cite{hld-ingestion}.

Thiết kế provider boundary làm cho hệ thống có thể chạy local không cần credential nhưng vẫn sẵn sàng chuyển sang provider thật. Nếu \texttt{LLM\_PROVIDER} không hợp lệ hoặc thiếu \texttt{OPENAI\_API\_KEY}, factory fallback về local parser. Nếu parser provider ném \texttt{ParserError}, worker đánh dấu parse job \texttt{failed} với \texttt{error\_code = llm\_parse\_failed} và bảo toàn file gốc.

\section{Embedding provider và semantic search}

Embedding provider boundary nằm tại \texttt{backend/src/jobconnect/integrations/embedding/}. \texttt{base.py} định nghĩa \texttt{EmbeddingProvider} protocol, \texttt{EmbeddingError} và constant \texttt{EMBEDDING\_DIM = 384}. \texttt{local.py} bọc \texttt{LocalHashEmbeddingProvider} dựa trên SHA-256-seeded bag-of-words embedder với \texttt{embedding\_version = "hash-v1"}. \texttt{openai.py} gọi endpoint \texttt{/v1/embeddings}, dùng \texttt{dimensions=384} và default model \texttt{text-embedding-3-small}; vector sai chiều phải raise \texttt{EmbeddingError}, không được silently padding \cite{hld-matching}.

\begin{table}[H]
\caption{Cấu hình embedding provider}
\centering
\begin{tabularx}{\textwidth}{p{0.33\textwidth}p{0.20\textwidth}X}
\toprule
\textbf{Biến môi trường} & \textbf{Default} & \textbf{Ý nghĩa}\\
\midrule
\texttt{EMBEDDING\_PROVIDER} & \texttt{local} & Chọn \texttt{local} hoặc \texttt{openai}; unknown fallback về local.\\
\texttt{OPENAI\_EMBEDDING\_API\_KEY} & unset & Credential riêng cho embedding; fallback sang \texttt{OPENAI\_API\_KEY}.\\
\texttt{OPENAI\_EMBEDDING\_MODEL} & \texttt{text-embedding-3-small} & Model name đưa vào \texttt{embedding\_version}.\\
\texttt{OPENAI\_EMBEDDING\_BASE\_URL} & OpenAI API & Cho endpoint compatible hoặc proxy.\\
\texttt{OPENAI\_EMBEDDING\_TIMEOUT\_SECONDS} & \texttt{30} & Timeout HTTP request.\\
\bottomrule
\end{tabularx}
\end{table}

Semantic search được tách khỏi matching. Job semantic search dùng query theo ý định ứng viên hoặc mô tả CV-style và rank theo JD requirement embeddings. Resume semantic search rank theo CV summary/experience embeddings, không rank title embeddings. Score semantic search là retrieval relevance, không phải final matching score \cite{requirements,hld-matching}. Sự tách biệt này ngăn người dùng hiểu nhầm rằng mọi kết quả semantic search đều là match đủ điều kiện nghiệp vụ.

\section{Matching pipeline}

Matching pipeline nhận một anchor đã đủ visibility: published job cho job-to-resume, active resume cho resume-to-job. Sau khi validate auth, role, ownership và anchor visibility, hệ thống lấy opposite pool, apply hard filters, score và optional rerank. Matching result không tạo application hoặc invite; nó chỉ là recommendation.

\subsection{Hard filters}

Hard filters đảm bảo cặp JD-CV có thể chấp nhận về mặt nghiệp vụ trước khi tính score semantic. Các rule gồm:

\begin{longtable}{p{0.28\textwidth}p{0.62\textwidth}}
\caption{Hard filters trong matching}\\
\toprule
\textbf{Filter} & \textbf{Rule}\\
\midrule
\endfirsthead
\toprule
\textbf{Filter} & \textbf{Rule}\\
\midrule
\endhead
\texttt{job\_type} & Job type phải match. Nếu JD là \texttt{remote}, location filter được bỏ qua.\\
\texttt{location} & Với job không remote, CV và JD phải có cùng location.\\
\texttt{seniority} & Seniority phải match chính xác.\\
\texttt{education} & CV education phải lớn hơn hoặc bằng job required education theo thứ tự \texttt{lop\_9 < lop\_12 < dai\_hoc < thac\_si < tien\_si}.\\
\texttt{certifications} & JD required certifications phải là subset của CV certifications.\\
Missing data & Thiếu dữ liệu hard-filter trên một phía làm pair fail hard filter.\\
\bottomrule
\end{longtable}

\subsection{Scoring deterministic}

Score formula giữ lại công thức prototype nhưng được đặt vào bối cảnh production MVP:

\begin{equation}
\begin{split}
final\_score = &\ 0.35 \times title\_sim + 0.35 \times skills\_sim \\
& + 0.20 \times req\_exp\_sim + 0.10 \times req\_summary\_sim \\
& + bonus\_exact\_skill - penalty\_missing\_required
\end{split}
\end{equation}

\begin{equation}
skills\_sim = 0.6 \times semantic\_skills + 0.4 \times exact\_overlap\_ratio
\end{equation}

Default \texttt{top\_k = 10}, allowed range \texttt{1..50}; default \texttt{min\_score = 0.7}, allowed range \texttt{0..1}. Bonus và penalty hiện giữ 0 đến khi có measured rule \cite{requirements,hld-matching}. Tie-break deterministic là \texttt{final\_score DESC}, sau đó \texttt{resume\_id ASC} cho job anchors hoặc \texttt{job\_id ASC} cho resume anchors. Quy tắc tie-break giúp cùng dữ liệu và cùng model version luôn trả thứ tự ổn định.

\subsection{Reranking và reasoning}

Runtime hiện tại cố gắng local cross-encoder rerank trên top deterministic candidates, dùng \texttt{sentence-transformers==3.0.1}. Nếu reranker lỗi hoặc unavailable, matching fallback về deterministic scoring và đưa warning trong runtime metadata \cite{implementation-matrix}. Đây là cách hợp lý cho MVP vì không làm matching fail chỉ vì model phụ không sẵn sàng.

Reasoning phải grounded in score breakdown và structured fields. Một reasoning tốt nên nói rõ kỹ năng nào overlap, hard filters nào pass, component nào mạnh/yếu và có thiếu embedding nào không. LLM-generated reasoning được phép, nhưng chỉ khi không invent fact. Điều này đặc biệt quan trọng trong tuyển dụng, vì một gợi ý sai hoặc không giải thích được có thể ảnh hưởng đến quyết định của người dùng.

\section{Authentication, authorization và auditability}

Authentication dùng email/password và JWT session. Implementation matrix ghi nhận custom HMAC JWT với \texttt{sub}, \texttt{role}, \texttt{iat}, \texttt{exp}; token hết hạn bị reject với \texttt{expired\_token}, token thiếu \texttt{exp} cũng bị reject \cite{implementation-matrix}. Disabled users không được publish, activate, apply, invite hoặc run matching.

Authorization theo role và ownership. Candidates quản lý profile, documents, resumes, applications, invites và notifications của mình. Recruiters quản lý jobs và recruiting actions liên quan đến job họ sở hữu hoặc được phép quản lý theo organization. Admin đọc dữ liệu monitoring, nhưng MVP không thêm destructive moderation. Privacy rule yêu cầu recruiters chỉ search/match active resumes, candidates chỉ search/match published jobs, và contact data của candidate phải hạn chế theo role/flow cần thiết \cite{hld-security}.

Audit logs là một phần quan trọng để hệ thống có thể vận hành. Các event cần audit gồm resume activated/archived, job published/closed, parse job failed, candidate applied, recruiter invite sent, invite accepted/rejected, application status changed và admin monitoring access where useful. Audit event cần actor user ID nếu có, target entity, event type, timestamp và metadata JSON.
